\section{The SAGA Comparison Theorem}\label{sec:5}

This section states and proves the central theorem of SAGA. Routine
verifications (such as componentwise checks of restriction naturality)
are compressed in the text, but every main argument that makes the
comparison theorem hold can be followed continuously within this
section.

\subsection{The central theorem}\label{sec:51}

Before the formal statement, we describe the shape of the claim in
plain terms. The stage is a selected monomorphic AAT cover, on which
stand the two complexes constructed independently in
Section~\ref{sec:4} --- the semantic side in the language of repair and
the equation side in the language of equations. The theorem asserts
that, under finitely many selection conditions matching the local data
of the two sides, (i) the two coefficient presheaves become naturally
isomorphic; (ii) this isomorphism identifies the two $H^1$ groups and
sends the obstruction class of the repair side to the obstruction class
of the equation side; and (iii) if moreover the family of repair states
satisfies the gluing condition (the true sheaf condition), then the
existence of a global repair is equivalent to the vanishing of the
obstruction class --- three stages in all. The inputs 1--8 below can be
read in four groups: inputs 1--2 are the two coefficients; inputs 3--4
are the correspondence connecting the coefficients and its
completeness; inputs 5--7 are the local state systems generating the
residuals and their correspondence; input 8 is the normalization of
empty overlaps.

\begin{theorem}[SAGA central theorem]\label{thm:central}
On a monomorphic AAT cover $\cU$, fix the following:
\begin{enumerate}
\item the abelian-group-valued presheaf $\Msem$ generated from
supported semantic atoms and a restriction-stable local repair relation
(\S\ref{sec:42});
\item the abelian-group-valued presheaf $\QE$ generated from an
architectural equation system $E$;
\item a restriction-natural map sending supported semantic atoms to
displayed Atom/equation residual classes;
\item completeness of the local repair relation with respect to that
map, and equation generator completeness (checked on each cover
intersection);
\item an affine semantic repair system $\Psem$ deriving an additive
torsor structure from the action of repair words, together with its
selected local repair atlas;
\item an equation-lift system $\PE$ acted on by $\QE$, together with
its selected local lift atlas;
\item a restriction-natural, generator-equivariant map
$\beta:\Psem\to\PE$ sending semantic local states to equation local
lifts (checked on each cover intersection);
\item for each empty cover intersection $V$ excluded from the products,
$\Msem(V)=\QE(V)=0$ and the subsingleton property of $\Psem(V)$ and
$\PE(V)$ (empty-overlap normalization).
\end{enumerate}
Then the following three stages hold.

\textbf{(i) SAGA Presentation Theorem.} The free extension
$\widetilde\chi_V:\Fsem(V)\to \QE(V)$ of the map of assumption~3
(\S\ref{sec:53}) satisfies, on each nonempty cover intersection $V$,
\[
\ker \widetilde\chi_V=\Rrep(V),
\qquad
\operatorname{im}\widetilde\chi_V=\QE(V).
\]
Of the kernel equality, the inclusion
$\Rrep(V)\subseteq\ker\widetilde\chi_V$ (repair-relation soundness) is
not an assumption: it is derived from the local-state data
(assumptions~5--7) in \S\ref{sec:53}. The reverse inclusion
$\ker\widetilde\chi_V\subseteq\Rrep(V)$ is repair-relation
completeness, and the image equality is equation-generator
completeness; both are assumption~4. By this exactness, the map induces
a natural isomorphism on the cover intersection diagram
\[
\Phi:\Msem\xrightarrow{\ \sim\ }\QE.
\]

\textbf{(ii) \v{C}ech comparison.} $\Phi$ induces a degreewise cochain
isomorphism between the \v{C}ech complexes built separately from
$\Msem$ and $\QE$,
\[
\kappa^\bullet:
C^\bullet_{\mathrm{sem}}(\mathcal U)
\xrightarrow{\ \sim\ }
C^\bullet_E(\mathcal U),
\]
and $\kappa$ commutes with the differentials. Hence an isomorphism
\[
\kappa_*:
H^1_{\mathrm{sem}}(\mathcal U)
\xrightarrow{\ \sim\ }
\check H^1(\mathcal U,\QE)
\]
is induced. For the residual $\rsem$ generated from the semantic local
repair atlas and the residual $\rE$ generated from the equation
local-lift atlas,
\[
\kappa_*([\rsem])=[\rE].
\]
Restricting the selected local atlases to each nonempty intersection
makes the local-state system on each intersection used by this
construction nonempty. The correspondence of residuals holds at the
cochain level, $\kappa^1(\rsem)=\rE$, when the two local atlases are
chosen aligned via $\beta$; when they are chosen independently, the
difference of the two cochains is an explicit $\delta^0$-image, and the
equality of classes holds.

\textbf{(iii) Grounded Global Gluing.} If moreover the $\Psem$ of
assumption~5 is a true semantic repair sheaf (defined in
\S\ref{sec:56}), then
\[
\mathrm{Nonempty}\,\Psem(W)
\iff
[\rsem]=0
\iff
[\rE]=0.
\]
\end{theorem}

\paragraph{The character of the three stages.}
In (ii), the cochain isomorphism and the induced $H^1$ isomorphism are
the general mechanism of \v{C}ech complexes that follows from a natural
isomorphism of coefficient presheaves. The latter half of (ii), the
residual correspondence, belongs to the proper content. The proper
content of this theorem is concentrated in (i), (iii), and the residual
correspondence: the independent construction of the semantic repair
presentation and the equation-generated quotient presentation
(Section~\ref{sec:4}); the construction of the generator map $\chi$
(Proposition~\ref{prop:chiE}); the derivation of relation soundness
from the local-state interpretation $\beta$ (\S\ref{sec:53}); the
placement of completeness / generation as conditions relativized to
architecture data; and the semantic correspondence of residual classes
(\S\ref{sec:55}).

\subsection{The structure of the proof}\label{sec:52}

The proof is organized in four stages, and
\S\ref{sec:53}--\S\ref{sec:56} give the arguments in this order:
\begin{center}
\small
\begin{tabular}{ll}
generator map $\chi$ $\to$ coefficient isomorphism
$\Phi : \Msem \simeq \QE$ & (\S\ref{sec:53})\\
$\to$ cochain isomorphism $\kappa$, $\kappa\delta=\delta\kappa$ &
(\S\ref{sec:54})\\
$\to$ $H^1$ isomorphism $\kappa_*$ & (\S\ref{sec:54})\\
$\to$ residual correspondence $\kappa_*([\rsem])=[\rE]$ &
(\S\ref{sec:55})\\
$\to$ grounded global gluing & (\S\ref{sec:56})\\
\end{tabular}
\end{center}

That soundness is a consequence rather than an assumption is the
structural crux of the theorem; its derivation opens \S\ref{sec:53}.

The comparison core (\S\ref{sec:53}--\S\ref{sec:55}) uses only local
data on cover intersections; the global sheaf condition does not appear
until the gluing of \S\ref{sec:56}. Which proof consumes which
assumption of Theorem~\ref{thm:central} is fixed as a table in
Appendix~\ref{app:A}.

The overall structure of the proof is shown in
Figure~\ref{fig:comparison}.

\begin{figure}[htbp]
\centering
\includegraphics[width=0.92\linewidth]{../zenodo_saga_figure1_comparison.png}
\caption{The structure of the SAGA comparison theorem
(Theorem~\ref{thm:central}). The semantic side (upper left) and the
equation side (upper right), constructed independently in
Section~\ref{sec:4}, pass through the coefficient isomorphism $\Phi$
(i) and the cochain isomorphism $\kappa$ commuting with the
differentials (ii) to the $H^1$ isomorphism $\kappa_*$ and the residual
class correspondence $\kappa_*([\rsem])=[\rE]$; under the true sheaf
condition the three-way equivalence for global repair (iii) holds.}
\label{fig:comparison}
\end{figure}

\subsection{Coefficient isomorphism: generation of soundness and presentation comparison}\label{sec:53}

The map of assumption~3 is a restriction-natural correspondence
\[
\chi_V:S(V)\longrightarrow \QE(V)
\]
sending supported semantic atoms to equation coefficients on each
nonempty cover intersection $V$ (the $\chi^E$ of
Proposition~\ref{prop:chiE} is a constructed instance of this input).
Since $\Fsem(V)$ is the free abelian group on $S(V)$ (\S\ref{sec:42}),
by the universal property $\chi_V$ extends to a unique homomorphism
$\widetilde\chi_V:\Fsem(V)\to \QE(V)$. The $\Phi$ to be constructed is
the map induced by $\widetilde\chi$ on the quotient by the repair
relation, and its well-definedness, injectivity, and surjectivity are
carried by three distinct conditions: soundness, completeness, and
generation. The joint validity of these three conditions is called
\textbf{SAGA presentation exactness}.

\paragraph{Generation of soundness.}
Take $x\in\Rrep(V)$ and $p\in\Psem(V)$ (under
Theorem~\ref{thm:central}, $\Psem(V)$ is nonempty by restriction of the
selected local atlas). By the action soundness of \S\ref{sec:42}, the
semantic action gives $x+p=p$. The generator-level equivariance of
$\beta$ extends to finite sums and additive inverses in the free group,
so
\[
\beta(p)=\beta(x+p)=\widetilde\chi(x)+\beta(p).
\]
Since the action of $\QE(V)$ on $\PE(V)$ is free, $\widetilde\chi(x)=0$.
Therefore
\[
\Rrep(V)\subset\ker\widetilde\chi_V.
\]
Deriving soundness from the local-state data instead of assuming it is
exactly this three-line argument.

The definition of $\beta$ in assumption~7 does not mention $\Rrep$;
all that is required is restriction naturality and generator-level
equivariance, so the derivation is not circular. That $\beta$ can be
constructed independently of soundness is shown by the concrete map in
the finite example of \S\ref{sec:58}.

\paragraph{Descent to the quotient.}
Since $\Rrep(V)\subset\ker\widetilde\chi_V$, the map
$\widetilde\chi_V$ factors through the quotient, inducing a unique
homomorphism
\[
\Phi_V:
\Msem(V)=\Fsem(V)/\Rrep(V)
\longrightarrow \QE(V).
\]

\paragraph{Injectivity.}
If $\Phi_V([x])=0$ then $x\in\ker\widetilde\chi_V$, and by
repair-relation completeness (assumption~4) $x\in\Rrep(V)$, hence
$[x]=0$.

\paragraph{Surjectivity.}
For any $q\in \QE(V)$, equation-generator completeness (assumption~4)
provides $x\in\Fsem(V)$ with $\widetilde\chi_V(x)=q$, and
$\Phi_V([x])=q$.

\paragraph{Naturality.}
On generators,
\[
\Phi_{V'}([\lambda|_{V'}])
  =\chi_{V'}(\lambda|_{V'})
  =\chi_V(\lambda)|_{V'}
  =\Phi_V([\lambda])|_{V'},
\]
where the second equality is the restriction naturality of $\chi$
(assumption~3). Since the generators generate $\Msem(V)$, $\Phi$
commutes with restriction on all elements.

Thus $\Phi:\Msem\simeq \QE$ is a natural isomorphism of presheaves on
the cover intersection diagram. This isomorphism is not a transport of
one complex from the other but a presentation comparison between the
semantic presentation and the equation quotient. The three conditions
carry three different mathematical jobs --- existence, injectivity, and
surjectivity of the map --- and removing any one of them destroys the
isomorphism, as counterexamples over the free group on one generator
$F=\Zint$ confirm:
\begin{itemize}
\item \emph{dropping soundness}: $R=2\Zint$, $Q=\Zint$, $\chi(n)=n$.
The relation $2=0$ does not become zero in $Q$, and the map does not
descend to the quotient.
\item \emph{dropping completeness}: $R=0$, $Q=\Zint/(2)$,
$\chi(n)=[n]$. The map descends and is surjective, but $2$ remains in
the kernel and the map is not injective.
\item \emph{dropping generation}: $R=2\Zint$, $Q=\Zint/(4)$,
$\chi(n)=[2n]$. The kernel coincides with $R$, but the image is
$\{0,[2]\}$ and the map is not surjective.
\end{itemize}

\subsection{Cochain commutation and the $H^1$ isomorphism}\label{sec:54}

\paragraph{Construction of the degreewise map.}
For each degree $n=0,1,2$, define
\[
\kappa^n:C^n_{\mathrm{sem}}(\mathcal U)\longrightarrow C^n_E(\mathcal U)
\]
by applying $\Phi$ to each intersection component:
\[
(\kappa^0a)_i=\Phi_{U_i}(a_i),\qquad
(\kappa^1c)_{ij}=\Phi_{U_{ij}}(c_{ij}),\qquad
(\kappa^2z)_{ijk}=\Phi_{U_{ijk}}(z_{ijk}).
\]
Since each $\Phi_V$ is an isomorphism, $\kappa^n$ is an isomorphism on
the product.

\paragraph{Commutation with the differentials.}
What must be shown is the commutativity of the diagram
\[
\begin{array}{ccccc}
C^0_{\mathrm{sem}}(\mathcal U)
&\xrightarrow{\ \delta^0_{\mathrm{sem}}\ }&
C^1_{\mathrm{sem}}(\mathcal U)
&\xrightarrow{\ \delta^1_{\mathrm{sem}}\ }&
C^2_{\mathrm{sem}}(\mathcal U)\\[2pt]
\big\downarrow{\scriptstyle\kappa^0}&&\big\downarrow{\scriptstyle\kappa^1}&&\big\downarrow{\scriptstyle\kappa^2}\\[2pt]
C^0_E(\mathcal U)
&\xrightarrow{\ \delta^0_E\ }&
C^1_E(\mathcal U)
&\xrightarrow{\ \delta^1_E\ }&
C^2_E(\mathcal U)
\end{array}
\]
For $a\in C^0_{\mathrm{sem}}(\cU)$ and $i<j$,
\[
\begin{aligned}
(\kappa^1\delta^0_{\mathrm{sem}}a)_{ij}
&=\Phi_{U_{ij}}(a_j|-a_i|)\\
&=\Phi_{U_j}(a_j)|-\Phi_{U_i}(a_i)|\\
&=(\delta_E^0\kappa^0a)_{ij},
\end{aligned}
\]
where the second line is the additivity and restriction naturality of
$\Phi$. Applying the same computation to the $(i,j,k)$ component of
$c\in C^1_{\mathrm{sem}}(\cU)$ (the alternating three-term sum) yields
$\kappa^2\delta^1_{\mathrm{sem}}=\delta^1_E\kappa^1$. Only the
homomorphy and naturality of $\Phi$ proved in \S\ref{sec:53} are used;
no new assumption enters here.

\paragraph{$H^1$ isomorphism.}
Since $\kappa$ is a degreewise isomorphism commuting with the
differentials, it preserves and reflects cocycles and coboundaries.
Hence
\[
\kappa_*:
H^1_{\mathrm{sem}}(\mathcal U)
\longrightarrow
\check H^1(\mathcal U,\QE),
\qquad
[c]\longmapsto[\kappa^1c]
\]
is well defined: replacing a representative by
$c+\delta^0_{\mathrm{sem}}a$ gives
$\kappa^1(c+\delta^0a)=\kappa^1c+\delta^0_E(\kappa^0a)$, and the class
does not change. The inverse map is induced in the same form by the
degreewise inverse $(\kappa^1)^{-1}$, and $\kappa_*$ is an isomorphism
of abelian groups.

\subsection{Residual correspondence}\label{sec:55}

Choose the semantic atlas $\{p_i\}$ and the equation atlas $\{e_i\}$
independently. On each chart, $\beta(p_i)$ and $e_i$ are elements of
the same $\QE(U_i)$-torsor $\PE(U_i)$, so there is a unique
$h_i\in \QE(U_i)$ with $e_i=h_i+\beta(p_i)$. On $U_{ij}$, using the
naturality and equivariance of $\beta$ (assumption~7),
\[
\begin{aligned}
e_j
&=h_j+\beta(p_j)\\
&=h_j+\beta(r_{\mathrm{sem},ij}+p_i)\\
&=h_j+\Phi(r_{\mathrm{sem},ij})+\beta(p_i)\\
&=(\Phi(r_{\mathrm{sem},ij})+h_j-h_i)+e_i.
\end{aligned}
\]
By uniqueness of differences in the equation torsor,
\[
r_{E,ij}=\Phi(r_{\mathrm{sem},ij})+h_j-h_i,
\qquad\text{that is,}\qquad
\rE=\kappa^1(\rsem)+\delta_E^0h.
\]
When the two atlases are chosen aligned via $\beta$, i.e.\ when
$e_i=\beta(p_i)$, we have $h=0$ and the cochain-level equality
$\kappa^1(\rsem)=\rE$ holds. When they are chosen independently, the
difference of the two cochains is the explicit coboundary
$\delta^0_Eh$, so the equality of classes
\[
\kappa_*([\rsem])=[\rE]
\]
holds. Since $\kappa_*$ is an isomorphism and preserves and reflects
zero, the equivalences $[\rsem]=0\iff[\rE]=0$ and
$[\rsem]\neq0\iff[\rE]\neq0$ follow as well.

This proves parts (i) and (ii) of Theorem~\ref{thm:central}
(\S\ref{sec:53}--\S\ref{sec:55}). The remaining three-way equivalence
(iii) is proved in the next subsection.

\subsection{Global repair: Grounded Global Gluing}\label{sec:56}

\begin{definitionstar}[true semantic repair sheaf]
A presheaf $\Psem$ on $S_X$ is called a \textbf{true semantic repair
sheaf} when: (1) $\Psem$ satisfies the sheaf condition for all covers
of the selected AAT topology; (2) the action of $\Msem$ commutes with
restriction and is locally free and transitive; (3) $\cU$ is a
monomorphic AAT cover belonging to that topology; (4) on the empty
intersections excluded from the products, $\Psem$ is subsingleton.
From these four conditions the sheaf condition for $\cU$ is derived,
and no per-cover amalgamation map is posited as separate data.
\end{definitionstar}

\begin{lemmafiveA}[ordered matching completion]
Suppose each morphism of the cover is a monomorphism and $\Psem$ is
subsingleton on the empty overlaps excluded from the products
(empty-overlap normalization). If a family $(q_i)$,
$q_i\in\Psem(U_i)$, agrees on all nonempty $i<j$ overlaps, then
$(q_i)$ is a matching family on all ordered overlaps.
\end{lemmafiveA}

\begin{proof}
Three kinds of overlap remain. On a self-overlap $U_{ii}$: since $u_i$
is a monomorphism, the diagonal $U_i\to U_i\times_W U_i$ is an
isomorphism, and under this canonical isomorphism the two projections
become the same morphism; hence the two restrictions of $q_i$ agree.
On a reversed overlap $U_{ji}$ ($i<j$): the pullback symmetry
$U_j\times_W U_i\cong U_i\times_W U_j$ is a canonical isomorphism
interchanging the two projections; restrictions are compatible along
this canonical isomorphism, so agreement on $U_{ji}$ is equivalent to
agreement on $U_{ij}$. On the empty overlaps excluded from the
products, $\Psem$ is subsingleton, so agreement is automatic.
\end{proof}

\begin{theorem}[Grounded Global Gluing]\label{thm:gluing}
In the setting of Theorem~\ref{thm:central}, if in addition $\Psem$ is
a true semantic repair sheaf, then the three-way equivalence
\[
\mathrm{Nonempty}\,\Psem(W)
\iff
[\rsem]=0
\iff
[\rE]=0
\]
holds. Moreover, for the correction $a$ obtained from a proof of
$[\rsem]=0$, there exists a unique global section whose restrictions
are the corrected family
\[
p_i^{\mathrm{corr}}=(-a_i)+p_i.
\]
\end{theorem}

\begin{proof}[Proof (forward)]
Suppose $[\rsem]=0$. Since the zero class is equivalent to the
existence of a matching correction (\S\ref{sec:42}), there is
$a\in C^0_{\mathrm{sem}}(\cU)$ with $\rsem=\delta^0_{\mathrm{sem}}a$.
The residual of the corrected family $p_i^{\mathrm{corr}}=(-a_i)+p_i$
is $\rsem-\delta^0a=0$ by choice independence (\S\ref{sec:42}), so the
corrected family agrees on all nonempty $i<j$ overlaps. By
Lemma~5.2A, the corrected family is a matching family on all ordered
overlaps. Since $\cU$ is a cover of the topology and $\Psem$ is a sheaf
for that topology, sheaf amalgamation provides a unique
$p\in\Psem(W)$ with $p|_{U_i}=p_i^{\mathrm{corr}}$.
\end{proof}

\begin{proof}[Proof (reverse)]
Conversely, suppose some $p\in\Psem(W)$ exists. On each chart, by
transitivity and freeness of the torsor there is a unique
$a_i\in\Msem(U_i)$ with $p_i=a_i+p|_{U_i}$. Taking differences on
overlaps,
\[
r_{\mathrm{sem},ij}=a_j-a_i,
\qquad\text{that is,}\qquad
\rsem=\delta^0_{\mathrm{sem}}a,
\]
so $[\rsem]=0$. The final equivalence of the three,
$[\rsem]=0\iff[\rE]=0$, is the residual class correspondence of
\S\ref{sec:55}.
\end{proof}

What is unique here is the amalgamation of the fixed corrected
matching family, not uniqueness of $\Psem(W)$ as a whole; the
difference between distinct global repairs is carried by degree-zero
data.

\paragraph{The equation side.}
The same conclusion descends to the equation side. Three additional
assumptions are made: SAGA presentation exactness holds on all
contexts of $S_X$; the $\beta$ of assumption~7 is given as a natural
transformation on the whole site (the comparison core uses only its
components on the intersection diagram, so this globalization is an
additional assumption); and $\PE$ is also a sheaf for the selected
topology. On each chart and nonempty overlap, the
$\Phi_V$-equivariant $\beta_V:\Psem(V)\to\PE(V)$ is a bijection:
injectivity follows from the transitivity of the semantic torsor, the
freeness of the target torsor, and the injectivity of $\Phi_V$;
surjectivity from the transitivity of the target torsor and the
surjectivity of $\Phi_V$. This local bijectivity lifts to the global
level in three steps. Injectivity: if $\beta_W(p)=\beta_W(p')$, then
on each chart the injectivity of $\beta_{U_i}$ gives
$p|_{U_i}=p'|_{U_i}$, and the separatedness of $\Psem$ gives $p=p'$.
Surjectivity: for $q\in\PE(W)$, set
$p_i:=\beta_{U_i}^{-1}(q|_{U_i})$ on each chart; by the naturality of
$\beta$ and injectivity on overlaps, the family $(p_i)$ agrees on all
nonempty $i<j$ overlaps, and by Lemma~5.2A and the amalgamation of
$\Psem$ it glues to a unique $p\in\Psem(W)$. Finally, $\beta_W(p)$
and $q$ have equal restrictions to each chart, so the separatedness
of $\PE$ gives $\beta_W(p)=q$. Thus
\[
\mathrm{Nonempty}\,\PE(W)\iff[\rE]=0
\]
follows.

\subsection{Scope of the theorem}\label{sec:57}

This comparison theorem is an additive theorem for
abelian-group-valued presheaves. The pointed-set-valued $H^1$ of
nonabelian torsors, higher coherence via 2-cocycles and gerbes, and
stack descent each require independent statements. None of these
conclusions is derived from the additive $H^1$ comparison.

Moreover, the $H^1$ treated by this theorem is the \v{C}ech
$H^1(\cU,-)$ relative to the selected monomorphic AAT cover $\cU$.
Identification with cover-independent sheaf cohomology requires the
additional conditions of refinement invariance or Leray-type
acyclicity, and this paper does not claim it.

Furthermore, the Lean formalization of this theorem is carried out
over a site whose context category is thin (the finite-meet poset
model of \S\ref{sec:33} is the representative example). Over general
$\ArchCtx(X)$, the gluing argument of \S\ref{sec:56}, including
Lemma~5.2A, uses monomorphisms and the universal property of
pullbacks, but that version has not been formalized (the details of
the hypotheses that thinness renders unnecessary are in
Appendix~\ref{app:A4}).

\subsection{Finite witness: independently generated circle comparison}\label{sec:58}

\begin{example}\label{ex:circle}
Take a monomorphic 4-cycle cover with four charts and four nonempty
overlaps $U_{01}, U_{12}, U_{23}, U_{03}$, and no nondegenerate triple
overlap.

\paragraph{Comparison of coefficients.}
On each nonempty intersection $V$, take the semantic support to be a
single generator $\sigma_V$ and the local repair relation to be
$2\sigma_V=0$ (writing $\Zint[\sigma_V]$ for the free abelian group on
the generator $\sigma_V$); the semantic side independently generates
\[
\Msem(V)=\Zint[\sigma_V]/(2\sigma_V)\cong\Ftwo.
\]
On the equation side, $\ObsE(V)=\Zint$ and $\IOb(V)=(2)$ give
$\QE(V)=\Zint/(2)$. For the interpretation $\chi_V(\sigma_V)=[1]$,
the three conditions of \S\ref{sec:53} can be checked directly:
\begin{center}
\small
\begin{tabular}{ll}
soundness: & $2\sigma_V \mapsto [2] = 0$\\
completeness: & $\ker(\Zint \to \Zint/(2)) = 2\Zint$\\
generation: & $[1]$ generates $\Zint/(2)$\\
\end{tabular}
\end{center}
Hence the construction of \S\ref{sec:53} gives $\Msem\simeq \QE$.
This is not a transport of one complex but a presentation comparison
between the semantic presentation $\Zint[\sigma]/(2\sigma)$ and the
equation quotient $\Zint/(2)$.

\paragraph{Local state system and affine transitions.}
On each chart and each nonempty overlap put $\Psem=\Ftwo$, with
$\Msem(V)\cong\Ftwo$ acting by addition. For an oriented edge $e$
(source chart $\to$ target chart), take the two restrictions to the
overlap to be
\[
\rho^{\mathrm{src}}_e(x)=x,
\qquad
\rho^{\mathrm{tgt}}_e(x)=x+t_e,
\]
where the transitions are
\[
(t_{01},t_{12},t_{23},t_{30})=(1,0,0,0).
\]
The last edge is $U_{03}$ oriented $3\to0$; in $\Ftwo$ the sign under
orientation reversal does not change anything. Since there is no
triple overlap, no additional functoriality condition constrains
these restrictions.

\paragraph{Computation of the residual.}
On each chart choose the local state $p_i=0$. On the overlap of an
edge $e$, the same $0$ is restricted from the source chart and from
the target chart, but the two restrictions are different maps. By the
definition of the residual in \S\ref{sec:42} (the difference of the
two restrictions on the overlap),
\[
r_{\mathrm{sem},e}
=\rho^{\mathrm{tgt}}_e(0)-\rho^{\mathrm{src}}_e(0)
=t_e.
\]
Even with all local states equal to $0$, the residual does not
vanish: what carries the mismatch is not the values of the states but
the difference of the two restrictions carrying the same $0$ to the
overlap --- the affine transition. Therefore
\[
\rsem=(1,0,0,0)\in\Ftwo^4.
\]

\paragraph{The equation side and the independent construction of $\beta$.}
Put $\PE=\Ftwo$ and give the affine transitions on oriented edges
independently by the same numbers $(1,0,0,0)$. The local-state
interpretation $\beta$ is the identity map on the $\Ftwo$ coordinates
on each chart and each nonempty overlap. This definition mentions
neither $\Phi$ nor the repair relation. Equivariance on generators
is, under $\chi(\sigma)=[1]$, the direct $\Ftwo$ computation
$\beta(\sigma+x)=x+1=\chi(\sigma)+\beta(x)$, and commutation with
restriction is a componentwise computation from the fact that the
transitions on both sides are given by the same numbers. Thus the
data of assumptions~5--7 are assembled independently of $\Rrep$, and
the derivation of \S\ref{sec:53} yields soundness
$\widetilde\chi(2\sigma)=[2]=0$ --- agreeing with the soundness row
checked directly in the coefficient comparison above. Choosing
$e_i=0$ on each chart --- which is the aligned choice
$e_i=\beta(p_i)$ of \S\ref{sec:55} --- gives, at the cochain level,
\[
\rE=\kappa^1(\rsem)=(1,0,0,0).
\]

\paragraph{Cocycle property and non-coboundary property.}
Since there is no nondegenerate triple overlap, $C^2(\cU)=0$ and
$\delta^1\rsem=0$ holds automatically. On the other hand, for any
$a\in C^0(\cU)$ the components of $\delta^0a$ are the
$a_j|-a_i|$ on the overlaps, and taking the edge sum around the
4-cycle, each chart's component appears exactly twice, so in $\Ftwo$
the sum is zero. The edge sum of $\rsem$ is $1+0+0+0=1$, so $\rsem$
is not a $\delta^0$-image. Therefore
\[
[\rsem]\ne0,\qquad
[\rE]\ne0,\qquad
\kappa_*([\rsem])=[\rE].
\]

\paragraph{Relation to real code.}
This example is a finite witness that the SAGA comparison preserves
not only zero classes but also nonzero classes. What erects the
nonzero class is not the number of charts but the
\textbf{cycle-without-a-face mechanism}: odd parity on a closed loop
plus the absence of a face (triple overlap) filling that loop. The
one-cent obstruction of Section~\ref{sec:7} is an instance of this
mechanism appearing in a real architecture (\S\ref{sec:71}).
\end{example}
